\documentclass[11pt]{article}

\usepackage[margin=1in]{geometry}
\usepackage{graphicx}
\usepackage{booktabs}
\usepackage{hyperref}

\title{Detecting Fraudulent Job Postings:\\Exploratory Analysis and Detection Strategy}
\author{}
\date{}

\begin{document}

\maketitle

\section{Problem Context}

Online job platforms face increasing risks from fraudulent job postings that exploit job seekers.
This analysis explores the EMSCAD dataset containing 17,880 job advertisements labeled as legitimate or fraudulent. The objective is to identify patterns associated with fraudulent postings and propose a practical detection strategy.

\section{Dataset Audit}

The dataset contains 17,880 job postings, of which 866 (4.8\%) are labeled as fraudulent.
This indicates a highly imbalanced classification problem.

The dataset contains a mixture of text fields (e.g., job descriptions and requirements),
structured job attributes (e.g., employment type and industry), and platform metadata
(e.g., presence of company logos and screening questions).

Several data quality considerations were observed:
\begin{itemize}
\item Binary variables were stored as text and required conversion.
\item Some fields such as salary\_range and benefits contain high levels of missing values.
\item Missing information may itself provide predictive signals for fraud detection.
\end{itemize}

\section{Key Insights from Exploratory Analysis}

\subsection{Company profile as a credibility signal}

Fraudulent postings frequently lack company profiles or contain extremely short descriptions of the employer. Legitimate postings typically provide richer company descriptions.

\subsection{Low-effort postings correlate with fraud}

Fraudulent postings are more likely to lack company logos and screening questions,
suggesting lower credibility of the posting.

\begin{figure}[h]
\centering
\includegraphics[width=0.6\textwidth]{figures/fraud_distribution.png}
\caption{Distribution of legitimate vs fraudulent postings}
\end{figure}

\subsection{Textual patterns}

Fraudulent postings often contain shorter requirements sections and vague job descriptions,
indicating potential linguistic signals useful for detection.

\section{Proposed Detection Approach}

A hybrid fraud detection approach combining textual features and structured metadata is recommended.

Key feature groups include:
\begin{itemize}
\item Employer credibility signals (company profile presence and length)
\item Text features derived from TF-IDF representations of job descriptions and requirements
\item Platform metadata such as company logos, screening questions, and telecommuting
\end{itemize}

A logistic regression model with TF-IDF features provides a strong and interpretable baseline,
with tree-based models (e.g., Random Forest or Gradient Boosting) as potential extensions.

\section{Evaluation and Deployment Strategy}

Given the class imbalance, evaluation should focus on:

\begin{itemize}
\item Precision
\item Recall
\item F1 score
\item Precision--Recall AUC
\end{itemize}

In practice, the model would function as a risk scoring system that flags suspicious postings for manual review.

\end{document}